\section*{Methods}

\subsection*{Problem formulation}

An item consists of a query $q$, evidence claims $E$, a query time and, where
relevant, an observer or possible-world identifier. A claim may be supported in
isolation but inapplicable to the query. We define
\begin{equation}
\mathrm{Apply}(e,q,S)\in\{\mathrm{valid},\mathrm{invalid},\mathrm{unresolved}\},
\end{equation}
where $S$ is the active situation state. The system first constructs or receives
$S$, then updates competing claims and selects an action from answering,
retrieval, clarification, abstention or task-specific action.

The end-to-end system is decomposed as
\begin{equation}
(q,E)\xrightarrow{f_{sense}}\hat S\xrightarrow{f_{update}}\hat S'
\xrightarrow{\pi}a\xrightarrow{g}y.
\end{equation}
This decomposition supports stage-specific evaluation. The gold-state condition
supplies $S$ directly and evaluates $f_{update}$, $\pi$ and the deterministic
answer renderer. A predicted-state condition applies a text-only deterministic
sensor before the same update and policy stages. It evaluates composition on
generated claim language, but not $f_{sense}$ on unrestricted natural language.

\subsection*{Situation schema}

The schema contains actors and entities; events and state transitions; valid
and observation time; scope; modality and source status; observer-specific
knowledge; possible-world identity; provenance; and unresolved answer-critical
variables. Update rules include temporal supersession, cancellation,
source-authority preference where specified, scope containment, observer
visibility and world separation. Unresolved variables are preserved rather than
converted into a guessed binary state.

\subsection*{SituationCatch-Bench generation}

The deterministic generator creates 600 instances for each of seven categories
using seeded lexical variation and shuffled claim order. Each item includes the
question, claims, structured situation fields, gold action, gold answer aliases
and category. The categories and gold semantics are documented in the bundled
datasheet. Generation and evaluation are intentionally co-designed: the
benchmark is a conformance suite for explicit semantics, analogous to unit tests
for a software protocol, rather than an independent estimate of natural-language
performance.

Future benchmark versions should use separate authors for data construction and
system development, clustered scenario families, hidden generators, natural and
semi-synthetic sources, adjudicated annotations, compositional splits and
adversarial state manipulations. Recommended natural-data fields and annotation
instructions are supplied in the Supplementary Information.

\subsection*{Reference implementation}

SCQA is a deterministic Python reference implementation with no trained
parameters. It parses the supplied structured fields, performs category-specific
state updates and selects a permitted response. The three software-control
baselines select claims using lexical overlap, recency or latest mention and do
not maintain the complete situation state. These controls test whether the
specified state variables are necessary for the generated invariants; they are
not intended to represent competitive LLM, RAG or agent baselines.

\subsection*{Predicted- and corrupted-state evaluation}

The predicted-state sensor receives only claim text and infers event type,
actor, temporal expression, modality, source status and world or scope using
documented deterministic patterns; it cannot copy hidden gold slots. The
downstream engine is unchanged. In a pre-specified corruption condition,
predicted source status is replaced by a generic reported source, modality is
forced to confirmed and scope is forced to global. All three conditions contain
the same 4,200 identifiers and are scored by exact action and answer. Item-level
and aggregate outputs are released in
\texttt{results/state\_conditions.csv} and
\texttt{results/state\_conditions\_summary.csv}.
The sensor is additionally scored claim by claim for actor, event, temporal
expression, validity interval, scope, modality, joint source/status,
observer-knowledge and world slots. Because claim order is preserved, this
comparison does not require learned alignment. The resulting files are
\texttt{results/predicted\_state\_slots.csv} and its aggregate summary.

\subsection*{Metrics and analysis}

The primary diagnostic measure is exact action-and-answer accuracy. For items
requiring clarification, correctness requires selecting clarification rather
than providing the latent answer. Category-level accuracy and matched component
ablations localize failure mechanisms. Because instances share template
families, item-level confidence intervals and $P$ values are not used as primary
scientific evidence in the revised manuscript. The released scripts retain
item-level summaries for regression testing, but future natural-language studies
should use scenario- or template-clustered bootstrap intervals, paired tests,
multiple-comparison control and pre-specified primary outcomes.

Rule-derived confidence values from the original implementation are not treated
as neural calibration. End-to-end work should report Brier score, expected
calibration error, selective risk, risk--coverage area, clarification utility
and abstention precision/recall using model- or state-derived probabilities.

\subsection*{Ablation design}

Seven ablations independently remove temporal validity, source status, modality,
scope, observer knowledge, world identity or missing-variable detection. Each
ablation is applied to all items, and its expected effect is specified before
execution. The purpose is to test causal alignment between representation and
failure category, not to estimate an average treatment effect in natural text.

\subsection*{Protocol for end-to-end multi-model evaluation}

A confirmatory study should compare at least the following matched conditions:
direct prompting; chain-of-thought or structured prompting; date-aware context;
standard RAG; temporal or event-aware RAG; reflection or verification agent; and
Situation Engineering. The generation model, evidence set, retrieval corpus,
maximum context, output budget and sampling parameters should be held constant
where the design permits. At least three proprietary and three open-weight model
families should be evaluated using immutable model identifiers.

The primary intervention contrast must be
\emph{same model + same evidence + same computation budget} with versus without
the explicit situation layer. A factorial analysis should independently vary
prompt, context/retrieval, harness and situation state. Competitive conditions
should include long-context prompting, few-shot structured reasoning,
self-consistency, self-refinement, Self-RAG, graph or temporal retrieval,
dialogue/belief-state tracking, clarification and calibrated abstention.

For each call, researchers should archive the prompt, evidence, raw response,
normalized answer, action decision, model identifier, provider, execution date,
latency, token use, monetary cost, errors and retries. Evaluation must report
gold-state, predicted-state and deliberately corrupted-state conditions. The
predicted-state condition should score entity/event extraction, temporal
relations, validity intervals, scope, modality, source status, observer
knowledge, possible-world identity and final answer. This enables error
attribution to sensing, updating, policy or generation.

The bundled \texttt{code/run\_multimodel\_eval.py} provides provider adapters and
refuses to create placeholder outputs when credentials are absent. No
frontier-model performance values are reported because such calls were not
executed in the study environment. The script is an executable protocol, not a
completed experiment.

\subsection*{Generalization and efficiency protocol}

External evaluation should pre-register held-out scenario families and test
unseen templates, longer event chains, nested scopes, conflicting authoritative
sources, multiple observers, multilingual questions and adversarial stale
evidence. Splits must occur at scenario or generator-grammar level rather than
over lexical variants. Uncertainty intervals should use scenario-clustered
resampling. Each condition must report latency, input and output tokens,
retrieval and tool calls, retries, monetary cost and clarification overhead.
Accuracy gains without matched resource reporting cannot identify the effect of
the situation layer.

\subsection*{Human and deployment evaluation}

Human evaluation should use at least three independent annotators per natural
item, blinded to method, with a documented adjudication procedure. Recommended
outcomes include state-slot agreement, answer correctness, appropriateness of
clarification or abstention, unnecessary clarification and perceived
justification quality. Studies involving participants should report ethics
approval or exemption, consent, compensation and privacy handling.

The released annotation program generates separately shuffled, blinded packets
for three or more annotators. It requires complete independent labels before
computing Fleiss' $\kappa$ and unanimous agreement for action, answer and state
slots; incomplete packets terminate with an error. Disagreements are adjudicated
only after independent files are frozen. No human agreement value is reported
because independent human annotations were not collected.

Deployment studies should additionally measure state drift, invalidation
latency, recovery after contradictory observations, action-precondition
violations, security attacks on state metadata, throughput, tail latency,
energy, token use and cost. The state and provenance required to audit a
consequential decision should be retained subject to privacy and access-control
requirements.

\subsection*{Reproducibility and reporting protocol}

The package contains the generator, generated benchmark, reference
implementation, predictions, figures, datasheet and exact commands. To make
future systems comparable, authors should report: (1) the situation schema;
(2) which variables are supplied, retrieved or inferred; (3) update and conflict
semantics; (4) response and action policies; (5) provenance and observability;
(6) stage-specific errors; (7) prompt, context, retrieval, memory, tool and
harness configurations; and (8) resource costs. A machine-readable experiment
manifest is recommended.

\subsection*{Ethics and use of generative AI}

The controlled benchmark contains no personal records or human participants.
Generative AI assisted drafting and software development. The author inspected
the code, generated outputs, calculations and references and remains responsible
for the manuscript. No empirical values were fabricated to represent unexecuted
model calls or human studies.
